\chapter{Problem Description and Motivation}

Route planning is a fundamental problem in transportation networks and algorithm design.
Classical navigation systems compute routes that minimize either travel distance or travel time.
However, in many real-world scenarios the primary cost of transportation is fuel consumption,
which depends not only on the length of the route but also on the price of fuel at the stations
the driver chooses to stop at.
On a long-distance Thai route such as Bangkok to Chiang Mai ($\approx$700\,km), a driver
with a 50-liter tank consuming 10\,km/L could spend between 1,750\,THB and 2,450\,THB
on fuel depending solely on which stations they choose---a difference of up to 40\% driven by
price variation alone.

Fuel prices vary by up to 3--5\,THB/L across Thai provinces (based on EPPO historical
data), due to supply chain differences, regional demand, and transportation costs.
This variation is large enough to make refueling strategy a first-class optimization concern,
yet existing navigation systems treat all fuel stops as equivalent.

This project studies the \textit{fuel-aware route optimization problem}.
Concretely, we want to answer: given a road network and a set of fuel stations with known
prices, what is the cheapest way to drive from a starting point to a destination without running
out of fuel?
The solution requires determining both:
\begin{itemize}
    \item the travel path between origin and destination, and
    \item the sequence of refueling decisions along that path (where to stop, and how much
          fuel to buy at each stop).
\end{itemize}

A critical subtlety that distinguishes this problem from standard shortest-path routing is the
\textit{variable purchase decision}: the driver can choose to buy any amount of fuel at each
station (subject to tank capacity), rather than always filling to capacity.
This seemingly small freedom significantly changes the algorithm design.
In the full-tank-only model, the fuel dimension disappears and standard Dijkstra on the road
graph suffices.
In the variable purchase model, the algorithm must track both the vehicle's position and its
current fuel level as a joint state, expanding the search space from $O(V)$ to $O(V \times F)$
where $F$ is the number of discretized fuel levels.
We elaborate on this in Section~\ref{sec:variable_purchase}.

Formally, given a road network $G = (V, E)$ and a set of fuel stations with associated prices
$p\colon V \to \mathbb{R}^{+}$, the goal is to compute a route and refueling strategy that
minimizes:
\[
    \min \sum_{i} b_i \cdot p(v_i)
    \quad\text{subject to: fuel level} \geq 0 \text{ at all times,}
\]
where $b_i$ is the amount of fuel purchased at station $v_i$ and the fuel-level constraint must
hold at every point along the route.

This problem combines concepts from graph algorithms, constrained optimization,
resource-bounded search, and transportation network analysis.
The primary aim of this project is empirical: we implement and benchmark six algorithmic
strategies on Thailand-scale road networks using real provincial price data, and propose
Partial-Fill A* (PF-A*) as a fuel-level-aware heuristic extension of the existing RF-A*
algorithm of Nandy et al.~\cite{nandy2023}.
To our knowledge, this is the first empirical study applying Gas Station Problem algorithms
to Thailand-scale road networks with real provincial price data.
